\section{Introduction}
\label{sec:intro}

\subsection{Motivation}
Degenerative cervical spine disease is among the most common reasons a cervical MRI is acquired.
Degenerative cervical myelopathy (DCM) is the leading cause of spinal-cord dysfunction in adults
worldwide, and cervical spondylotic radiculopathy is a frequent cause of arm pain and disability. The
radiologist's report on a cervical MRI is built on a set of quantitative measurements: the heights and
shape of each vertebral body, the dimensions and signal of each intervertebral disc, the antero-posterior
diameter of the spinal canal and the cord and the space available between them, the sagittal alignment
(lordosis), and the presence of abnormal cord signal. These numbers are then compared, often
implicitly, against normative values to decide whether a finding is present and how severe it is.

Performed manually, this is slow and variable. Measuring a single disc height, a canal diameter, or a
Cobb angle by hand requires careful caliper placement on the right slice, and studies of inter-observer
agreement repeatedly show that different readers --- and the same reader on different days --- produce
materially different numbers. The non-AI baseline our application competes against is therefore a
manual workflow that is both labour-intensive and noisy.

\subsection{What this system does}
MRI-ReportGenerator takes one sagittal T2 cervical MRI and produces a structured report. The pipeline is:
\begin{center}
\texttt{Input $\rightarrow$ Segmentation $\rightarrow$ Measurements $\rightarrow$ Interpretation $\rightarrow$ Report}
\end{center}
Segmentation is delegated to three established, openly licensed deep-learning tools, each chosen for the
anatomical structure it segments best. The measurement layer is organized into five groups that mirror
the structures a radiologist reports on. The interpretation layer compares each measurement to a central
catalog of thresholds, each carrying an explicit literature citation and a modality caveat. Every output
is phrased as a finding ``flagged for physician review''; the system never asserts a diagnosis.

\subsection{Positioning: an application, and an honest one}
This project is positioned as an \emph{application}, not a research contribution claiming a new
state-of-the-art algorithm. Its value is in the disciplined engineering: composing validated tools,
grounding every clinical number in a primary source, and --- crucially --- validating the assembled
pipeline against real healthy and pathological data with a methodology appropriate to the available
ground truth. Two principles run through the work:
\begin{enumerate}[leftmargin=*]
  \item \textbf{Measurements are disease-agnostic geometry/signal detectors anchored on healthy norms.}
  A canal narrowed to \SI{8}{\milli\meter} reads as stenosis whether the cause is spondylosis, tumour,
  or trauma. The normal range for each measurement is taken from healthy-population literature, not
  learned from a disease cohort. Consequently the system's abnormal-detection generalizes across
  pathologies that produce the same anatomical change, and we never overfit to a particular disease.
  \item \textbf{Honesty about uncertainty.} Where a threshold is borrowed from a different modality
  (e.g.\ a radiograph-derived ratio applied to MRI), or where no cited value exists, the system says so
  explicitly rather than inventing a number.
\end{enumerate}

\subsection{The validation problem}
The central difficulty is that \emph{no public dataset pairs cervical MRI images with per-case expert
measurements}. We confirmed this across repeated, exhaustive dataset hunts. Without such ground truth we
cannot compute a true sensitivity or specificity against a radiologist. We therefore adopt a
\textbf{threshold-crossing and distribution-separation} design: on healthy necks every measurement must
sit on the normal side of its cited threshold and never cross; on a symptomatic cohort the distribution
must shift into the abnormal side, with visibly pathological cases flagging. This demonstrates that each
measurement ``reads the right range'' without claiming per-case diagnostic accuracy.

\subsection{Contributions}
\begin{itemize}[leftmargin=*]
  \item An end-to-end, deployable cervical-MRI measurement and reporting pipeline composing three
  segmentation engines, five measurement groups, and a cited-threshold interpretation layer.
  \item A central, fully cited threshold catalog with explicit modality caveats and a no-diagnosis
  output ontology.
  \item A literature-grounded methodology for the two hardest measurements (the vertebral
  endplate/corner geometry and the cervical Cobb angle), including the discovery that fitting a line to
  SPINEPS' own predicted endplate voxels outperforms both corner-extrema and corpus-border methods.
  \item A reproducible threshold-crossing validation on real cervical MRIs with non-parametric statistics,
  showing strong separation for the canal/cord group, a correctly null result for the vertebral
  compression screen on a non-compression cohort, the detection of a real sign-level bug in the
  disc-height index, and --- on scaled, balanced cohorts --- the honest revision of an initial alignment
  lead to a \emph{non-discriminator} (and the disc signal/bulge axes to negatives).
  \item A transparent ``development journey'' documenting every substantive mistake, how it was found,
  the fix, and how the fix was validated --- the methodological core of the project.
\end{itemize}

\subsection{Paper organization}
Section~\ref{sec:clinical} reviews the clinical background and the measurements.
Section~\ref{sec:related} surveys the segmentation tools, prior measurement methods, and datasets.
Section~\ref{sec:arch} describes the system architecture. Section~\ref{sec:data} details the datasets.
Sections~\ref{sec:seg}--\ref{sec:g6} present the segmentation and the five measurement groups plus the
interpretation layer. Section~\ref{sec:valmethod} formalizes the validation methodology and
Section~\ref{sec:results} reports the results. Section~\ref{sec:journey} narrates the development
journey, Section~\ref{sec:limits} the limitations, Section~\ref{sec:deploy} the deployment, and
Section~\ref{sec:conclusion} concludes.
